\documentclass[11pt,a4paper]{article}
\usepackage[T1]{fontenc}
\usepackage{isabelle,isabellesym}
\usepackage{amssymb}
\usepackage{amsmath}
\usepackage[only,bigsqcap,bigparallel,fatsemi,interleave,sslash]{stmaryrd}
\usepackage{pdfsetup}

\urlstyle{rm}
\isabellestyle{it}

\begin{document}

\title{Napoleon's Theorem in Isabelle/HOL}
\author{Arthur Freitas Ramos \and David Barros Hulak \and Ruy J. G. B. de Queiroz}
\maketitle

\begin{abstract}
We formalize Napoleon's theorem in Isabelle/HOL.  Given a nondegenerate
Euclidean triangle, equilateral triangles constructed externally on its three
sides have centres that form an equilateral triangle.  The construction is
represented in complex coordinates, and both orientations of the equilateral
rotation are covered by one parametrized development.
\end{abstract}

\tableofcontents

\section{Introduction}

Napoleon's theorem is a classical triangle-geometry result
\cite{honsberger1995episodes}: the three
equilateral triangles erected on the sides of an arbitrary triangle determine
three centres which themselves are the vertices of an equilateral triangle.
The orientation of the erected triangles matters.  A coordinate
formalization must therefore distinguish the two rotations through
\(\pm\pi/3\), and it must make explicit which orientation gives the external
construction.

The exterior condition is part of the formal specification.  For a directed
side \(ab\), the signed side orientation of a point \(x\) is
\(\operatorname{Im}((x-a)/(b-a))\).  An apex is external precisely when its
side orientation has strictly opposite sign to that of the third vertex.

The formalization uses the complex plane.  If \(s\) is either
\(\operatorname{cis}(\pi/3)\) or \(\operatorname{cis}(-\pi/3)\), the apex on the
directed side from \(x\) to \(y\) is
\[
  E_s(x,y)=x+s(y-x).
\]
The centre of the corresponding equilateral triangle is the average of its
three vertices.  The published AFP development of Morley's theorem
\cite{morleyAFP} supplies
the equilateral-triangle characterization used by the present entry; the
present theory adds the cyclic centre calculation specific to Napoleon's
configuration.

\section{Definitions and construction correctness}

The theory defines the oriented apex and the centre of an equilateral triangle
on a directed side.  For distinct endpoints, the rotation characterization
immediately proves the two equal-length equations that certify the
construction as equilateral.  Keeping the rotation parameter explicit avoids
duplicating the proof for clockwise and counterclockwise constructions.
The predicate \texttt{external\_side} states that two points lie strictly on
opposite sides of the directed line.  The theorem
\texttt{napoleon\_external\_construction} proves this predicate for all three
cyclic sides, in addition to the six equal-distance facts for the constructed
equilateral triangles.

\section{The centre calculation}

Let \(N_{AB}\), \(N_{BC}\), and \(N_{CA}\) denote the three centres.  The
complex rotation parameter satisfies
\[
  s^2=s-1,
\]
and the cyclic differences simplify to
\[
  N_{CA}-N_{AB}=(1-s)(N_{BC}-N_{AB}).
\]
Since \(1-s\) is the opposite \(\pm\pi/3\) rotation, the two vectors from
\(N_{AB}\) to the other two centres have equal lengths and meet at an
equilateral angle.  The final theorem is stated through the corresponding
pairwise-distance characterization.  The final bundled corollary
\texttt{napoleon\_external\_theorem} packages this centre conclusion together
with all three external equilateral constructions under the non-collinearity
assumption, so the formal statement mirrors the complete English theorem.

\section{Formalization notes}

The entry imports the released \texttt{Morley\_Theorem} session rather than
duplicating its complex equilateral infrastructure.  The \texttt{napoleon\_centre}
is the arithmetic mean of the three vertices, hence the centroid; on an
equilateral triangle this agrees with the circumcenter, incenter, and
orthocenter.  The source theory is kept as a single top-level AFP theory, with
the document containing the mathematical explanation and proof sketch.

\input{session}

\bibliographystyle{abbrv}
\bibliography{root}

\end{document}
